\subsection{Conjugation in the affine group}

From \eqref{eq:affine-group-law},
\begin{align*}
 (q,v)^{-1}(p,u)(q,v)
 &=(q^{-1},-q^{-1}v)(p,u)(q,v)\\
 &=\left(p,q^{-1}(u-v)+q^{-1}pv\right)\\
 &=\left(p,q^{-1}\bigl(u+(p-1)v\bigr)\right),
\end{align*}
which proves \eqref{eq:affine-conjugation}.  The multiplier of the left input is
preserved by conjugation.  Consequently each fixed-multiplier set is a subquandle,
and taking $p=q=t$ gives \eqref{eq:alexander-quandle-operation}.

The torsion expression can be recognized as the difference between the translation
parts of the two affine products.  More precisely, define the translation-coordinate
projection by $\operatorname{trans}(p,u)=u$.  Then
\[
 \operatorname{trans}\bigl((p,u)(q,v)\bigr)
 -\operatorname{trans}\bigl((q,v)(p,u)\bigr)
 =(1-q)u+(p-1)v.
\]
This verifies \eqref{eq:affine-torsion} and its fixed-branch restriction.

\subsection{Twisted cocycle and primitive}

Write $c=t-1$, $\alpha=t^{-1}$, and let the coefficient action be multiplication
by $\beta$.  For $\phi(x,y)=c(y-x)$, the left-hand side of
\eqref{eq:twisted-cocycle-equation}, divided by $c$, is
\begin{align*}
 \beta(y-x)+z-\bigl(\alpha x+(1-\alpha)y\bigr)
 =z-(\alpha+\beta)x+(\alpha+\beta-1)y.
\end{align*}
The right-hand side is
\begin{align*}
 &\beta(z-x)+(1-\beta)(z-y)\\
 &\quad+\bigl(\alpha y+(1-\alpha)z\bigr)
       -\bigl(\alpha x+(1-\alpha)z\bigr)\\
 &=z-(\alpha+\beta)x+(\alpha+\beta-1)y.
\end{align*}
Thus $\phi$ is a twisted cocycle for every $\beta$.

For $f(x)=Cx$, convention \eqref{eq:twisted-coboundary-convention} gives
\[
 \delta_T^+f(x,y)=C(\alpha-\beta)(y-x).
\]
If $\alpha-\beta$ is invertible, choosing $C=c/(\alpha-\beta)$ proves
nonresonant exactness.  At $\beta=\alpha$, this calculation vanishes for every
linear $f$, but Theorem~\ref{thm:resonant-affine-class} is stronger: it excludes
every set-theoretic one-cochain over the finite field.

For completeness, the induction used there follows from
$g(\alpha x)=\alpha g(x)+cx$:
\begin{align*}
 g(\alpha^{j+1}x)
 &=\alpha g(\alpha^jx)+c\alpha^jx\\
 &=\alpha^{j+1}g(x)+(j+1)c\alpha^jx.
\end{align*}

\subsection{The extension calculation}

Relative to the ordered basis of $E=N\oplus X_t$, the operator in
\eqref{eq:extension-action} is
\[
 [T_E]=
 \begin{pmatrix}
 \alpha&-(t-1)\\
 0&\alpha
 \end{pmatrix},
 \qquad
 \det[T_E]=\alpha^2\ne0.
\]
Thus it defines an Alexander module.  With $s(x)=(0,x)$,
\begin{align*}
 T_Es(x)&=(-(t-1)x,\alpha x),\\
 (1-T_E)s(y)&=((t-1)y,(1-\alpha)y).
\end{align*}
Their sum is $(\kappa_t(x,y),x\triangleright y)$, which proves
\eqref{eq:kappa-obstruction}.  Equivalently, the Alexander operation on $E$ is
\[
 (a,x)*(b,y)
 =\bigl(\alpha a+(1-\alpha)b+\kappa_t(x,y),
         x\triangleright y\bigr).
\]
This is the explicit extension behind
Theorem~\ref{thm:resonant-state-sum-collapse}.

\subsection{Two coloring-count checks}

The state sum uses crossing weights modified by the Alexander numbering of the
source region.  The extension theorem makes the total weight zero for every planar
coloring, so only the number of colorings remains.  Two small cases check the
result.

For the positive two-strand braid, use the coloring map
\[
 B(a,b)=(b,a\triangleright b).
\]
Writing $d=b-a$, iteration gives $d_{j+1}=-\alpha d_j$.

The closure of $B^2$ is the positive Hopf link.  Its fixed-point equation forces
$(1-\alpha)d=0$, so $d=0$ because $t\ne1$.  Hence
\begin{equation}
 \Phi_{\kappa_t}(H)=q[0].
 \label{eq:hopf-collapsed-state-sum}
\end{equation}

The closure of $B^3$ is the positive trefoil.  Its fixed-point equation is
$(\alpha^2-\alpha+1)d=0$, giving
\begin{equation}
 \Phi_{\kappa_t}(3_1)=
 \begin{cases}
  q^2[0],&\alpha^2-\alpha+1=0,\\
  q[0],&\alpha^2-\alpha+1\ne0.
 \end{cases}
 \label{eq:trefoil-collapsed-state-sum}
\end{equation}
For $q=3$ and $t=\alpha=2$, the values are $3[0]$ and $9[0]$.  Their difference
is entirely the Alexander-quandle coloring count, not a cocycle enhancement.

\subsection{Variable-multiplier defect}

Let $K=\kappa(x,y)$.  From \eqref{eq:affine-conjugation},
\begin{align*}
 \kappa(x\triangleright y,z)-\kappa(x,z)
 &=\frac{1-r}{q}K,\\
 \kappa(x\triangleright z,y\triangleright z)
 &=\frac1rK.
\end{align*}
Therefore
\begin{align*}
 \mathfrak A_\kappa(x,y,z)
 &=\left(1+\frac{1-r}{q}-\frac1r\right)K\\
 &=\frac{(q-r)(r-1)}{qr}K,
\end{align*}
which proves \eqref{eq:variable-multiplier-anomaly}.  The translation coordinate
$w$ of $z$ cancels identically; this does not provide the higher coherence absent
from the construction.

\subsection{The figure-eight word}

With $a=(t,0)$, $b=(1,1)$, $A=a^{-1}$, and $B=b^{-1}$, multiply the letters of
$abbbaBAAB$ from left to right using \eqref{eq:affine-group-law}:
\begin{center}
\small
\begin{tabular}{ccc}
\toprule
Prefix & Multiplier & Translation \\
\midrule
$a$ & $t$ & $0$ \\
$ab$ & $t$ & $t$ \\
$abb$ & $t$ & $2t$ \\
$abbb$ & $t$ & $3t$ \\
$abbba$ & $t^2$ & $3t$ \\
$abbbaB$ & $t^2$ & $3t-t^2$ \\
$abbbaBA$ & $t$ & $3t-t^2$ \\
$abbbaBAA$ & $1$ & $3t-t^2$ \\
$abbbaBAAB$ & $1$ & $3t-t^2-1$ \\
\bottomrule
\end{tabular}
\end{center}
The final translation is $-(t^2-3t+1)$, proving
\eqref{eq:figure-eight-free-word} under the declared convention.
